\section*{अध्याय \ref{ch:b2:waves-on-strings} --- तारों और छड़ों पर तरंगें: दालांबेर समीकरण}

\begin{solution}{exo:b2:waves-on-strings:1}
$c = 2Lf = \qty{143}{m/s}$; $T = \mu c^2 = \qty{61}{N}$; तार पर $\lambda = 2L = \qty{1.3}{m}$,
हवा में $340/110 = \qty{3.1}{m}$; पाँचवाँ परदा: $110 \times 2^{5/12} =
\qty{147}{Hz}$।
\end{solution}

\begin{solution}{exo:b2:waves-on-strings:2}
$\sqrt{E/\rho}$: इस्पात \qty{5.1}{km/s}, ऐलुमिनियम \qty{5.1}{km/s}, सीसा
\qty{1.2}{km/s}। \qty{1}{km}: पटरी में \qty{0.20}{s}, हवा में \qty{2.9}{s} ---
अर्थात् दो ध्वनियाँ, पहले पटरी वाली, और बीच में \qty{2.7}{s} का अंतर।
\end{solution}

\begin{solution}{exo:b2:waves-on-strings:3}
$400$, $600$, $800$ Hz, और भीतरी निस्पंद $1$, $2$, $3$; विधा 3 के निस्पंद
$L/3$ और $2L/3$ पर। $T \times 2$: $200\sqrt2 = \qty{283}{Hz}$; $L/2$: \qty{400}{Hz}; $\mu
\times 3$: $200/\sqrt3 = \qty{115}{Hz}$।
\end{solution}

\begin{solution}{exo:b2:waves-on-strings:4}
$c = \sqrt{45/0.2} = \qty{15}{m/s}$; \qty{1.6}{s} में \qty{24}{m}; स्पंद उलटा
लौटता है ($r = -1$); और मुक्त छल्ले के साथ सीधा ($r = +1$)।
\end{solution}

\begin{solution}{exo:b2:waves-on-strings:5}
(क) $\mu = \rho\pi d^2/4 = \qty{6.1}{g/m}$; $c = 2Lf = \qty{334}{m/s}$; $T = \mu c^2 =
\qty{690}{N}$। (ख) $\sigma = T/A = \rho c^2 = \qty{0.87}{GPa}$, अवकाश $2.3$। (ग) $T/4 =
\qty{170}{N}$, और प्रतिबल वही (वह केवल $\rho c^2$ पर निर्भर है)। (घ) लंबा
तार पेटी में समाता नहीं; मोटा सादा तार बहुत कड़ा (और तीव्रता से
असंनादी) होता है और सेतु पर मुड़ता नहीं; जबकि ताँबे की लपेट
बिना कड़ेपन के द्रव्यमान जोड़ देती है।
\end{solution}

\begin{solution}{exo:b2:waves-on-strings:6}
(क) $m_n = 4L^2f^2\mu/n^2g = 6.6/n^2$ kg: $6.6$, $1.65$, $0.73$, \qty{0.41}{kg}। (ख)
$n^2 = 13.2$, $n = 3.6$: अनुनाद नहीं, वह विधा 3 और 4 के बीच है (निकटतम: 4,
\qty{0.41}{kg})। (ग) व्यावहारिक रूप से निस्पंद: चालक का आयाम फंदों के
सामने नन्हा होता है। (घ) हर आवर्तकाल में ऊर्जा कला में भरती जाती है और
जमा होती है; और अवमंदन (हवा, आंतरिक घर्षण, घिरनी से रिसाव) आयाम को
सीमित कर देता है।
\end{solution}

\begin{solution}{exo:b2:waves-on-strings:7}
$c = \qty{100}{m/s}$, $k = \qty{6.3}{rad/m}$, $\lambda = \qty{1.0}{m}$; $v_{\max} = \omega A =
\qty{1.3}{m/s}$, $a_{\max} = \omega^2A = \qty{790}{m/s^2}$; ढाल $kA = 0.013$; $Z =
\sqrt{T\mu} = \qty{1.0}{kg/s}$; $\langle\mathcal P\rangle = \tfrac12Z\omega^2A^2 = \qty{0.79}{W}$;
प्रति तरंगदैर्घ्य $\mathcal P/f = \qty{7.9}{mJ}$।
\end{solution}

\begin{solution}{exo:b2:waves-on-strings:8}
(क) $c_1 = \qty{100}{m/s}$, $c_2 = \qty{50}{m/s}$; $Z_1 = \qty{0.20}{kg/s}$, $Z_2 = \qty{0.40}{kg/s}$।
(ख) हलके तार से: $r = -1/3$, $t = 2/3$; भारी तार से: $r =
+1/3$, $t = 4/3$। (ग) दोनों दशाओं में $R = 1/9$, $\mathcal T = 4 \times 0.08/0.36 = 8/9$।
(घ) पार गया: \qty{0.67}{cm} ऊँचा, \qty{5}{cm} चौड़ा (वही अवधि, आधी
चाल); परावर्तित: \qty{-0.33}{cm}, \qty{10}{cm}।
\end{solution}

\begin{solution}{exo:b2:waves-on-strings:9}
(क) $E = F\ell/A\Delta\ell = 80 \times 2/(7.85 \times 10^{-7} \times 10^{-3}) = \qty{2.0e11}{Pa}$।
(ख) $c = \qty{5.1}{km/s}$। (ग) अनुदैर्घ्य: $c/2\ell = \qty{1.3}{kHz}$; अनुप्रस्थ:
$\sqrt{T/\mu} = \qty{114}{m/s}$, $f = \qty{29}{Hz}$ --- तनाव तार को
\qty{100}{MPa} तक ही प्रतिबलित करता है, अर्थात् $5 \times 10^{-4}$ गुना $E$, इसलिए अनुप्रस्थ चाल
अनुदैर्घ्य चाल की $\sqrt{5 \times 10^{-4}}$ गुनी है। (घ) $k = Ea = \qty{50}{N/m}$; $m
= \rho a^3 = \qty{1.2e-25}{kg}$ (लगभग $70$ परमाणु द्रव्यमान मात्रक)।
\end{solution}

\begin{solution}{exo:b2:waves-on-strings:10}
(क) $b_n = \frac2L\int_0^Ly(x, 0)\sin(n\pi x/L)\dd x$, जहाँ त्रिभुज $y =
2hx/L$ है $[0, L/2]$ पर (और सममित): खंडशः समाकलन कीजिए, $b_n = 8h\sin(n\pi/2)
/n^2\pi^2$। (ख) सम $n$ लुप्त हो जाते हैं: क्योंकि जिस मध्य-बिंदु को खींचा गया वह
हर सम विधा का निस्पंद है। (ग) $E_n = \tfrac14\mu L\omega_n^2b_n^2 \propto n^2/n^4 = 1/n^2$
(विषम $n$): अतः मूल विधा के पास $1/\sum_{\text{विषम}}n^{-2} = 8/\pi^2 = 81\%$ है। (घ)
$b_n \propto \sin(n\pi/5)/n^2$: अतः $n = 5, 10, \dots$ लुप्त हो जाते हैं। सेतु के पास
गुणक $\sin(n\pi x_0/L)$ नीची विधाओं के लिए छोटा रहता है और ऊँची विधाओं के लिए
बढ़ता जाता है: अतः अपेक्षाकृत अधिक ऊँचे संनादी, और चमकीली ध्वनि।
\end{solution}

\begin{solution}{exo:b2:waves-on-strings:11}
(क) $-m\omega^2 = k(2\cos ka - 2) = -4k\sin^2(ka/2)$। (ख) $v_\varphi = 2\omega_0\sin(ka/2)
/k$, $v_g = \omega_0a\cos(ka/2)$; दोनों $\to \omega_0a = c$ जब $ka \to 0$; और $ka = \pi$ पर:
$v_\varphi = 2c/\pi$, $v_g = 0$। (ग) पड़ोसी विपरीत कला में चलते हैं: अर्थात् कोई \omterm{prop:b2:waves-on-strings:modes}{अप्रगामी
तरंग} जो कुछ नहीं ढोती; और कोई $k$, जो $\pi/a$ से बड़ा हो, वही विस्थापन-समुच्चय देता है
जो कोई $k' = k - 2\pi/a$ देता --- कोई नई तरंग नहीं। (घ) $\omega_0 = c/a =
\qty{2.0e13}{rad/s}$, $f_{\max} = \omega_0/\pi = \qty{6.5}{THz}$; \qty{1}{MHz} पर: $\lambda =
\qty{5.1}{mm} = 2 \times 10^7a$, $ka \sim 3 \times 10^{-7}$: अर्थात् पूर्णतः अपरिक्षेपी।
\end{solution}

\begin{solution}{exo:b2:waves-on-strings:12}
(क) $T(x) = \mu gx$, $c = \sqrt{gx}$। (ख) $t = \int_0^L\dd x/\sqrt{gx} = 2\sqrt{L/g} =
\qty{2.0}{s}$; और स्पंद का ऊपरी भाग निचले से तेज़ चलता है: अतः वह
खिंच जाता है। (ग) $c = \sqrt{T/\mu}$ नोक की ओर बढ़ता है। (घ) ऊर्जा $\approx
\mathcal P\tau = ZV^2\tau = \sqrt{T\mu}V^2\tau$ अचर: अतः $V \propto \mu^{-1/4}$; और
$10^4$ का गुणक, $\mu$ में, $\times10$ देता है: \qty{100}{m/s} --- असली चाबुक
बाक़ी उस खुलती कुंडली से पाता है, जो ऊर्जा को और संकेंद्रित कर देती है,
और नोक ध्वनि की चाल पार कर जाती है।
\end{solution}

\begin{solution}{pb:b2:waves-on-strings:1}
\textbf{1.} $\mu = \rho\pi d^2/4 = \qty{6.1}{g/m}$; $c = 2Lf = \qty{334}{m/s}$; $T = \mu c^2
= \qty{685}{N}$।

\textbf{2.} $\sigma = \rho c^2 = \qty{0.87}{GPa}$: अवकाश $2.3$। $T \propto f^2$: अतः
आधे स्वर की तीव्रता $12\%$ जोड़ देती है, और एक अष्टक ऊपर उसे चौगुना कर देती है,
\qty{3.5}{GPa} --- तब वह टूट जाता है।

\textbf{3.} $230 \times 685 \approx \qty{160}{kN}$, अर्थात् सोलह टन।

\textbf{4.} $c = 2Lf = \qty{105}{m/s}$; \qty{685}{N} पर: $\mu = T/c^2 = \qty{63}{g/m}$,
अर्थात् \qty{3.2}{mm} की इस्पात-छड़ --- जो संनादी रूप से कंपने या सेतु पर
मुड़ने के लिए कहीं अधिक कड़ी है; जबकि ताँबे से लपेटा पतला क्रोड
बिना कड़ेपन के द्रव्यमान दे देता है।

\textbf{5.} $440$, $880$, $1320$, $1760$, $2200$, $2640$, $3080$, \qty{3520}{Hz}: $2$,
$4$, $8$ अष्टक हैं; $3$, $6$ किसी अष्टक से ऊपर पंचम; $5$ दो अष्टकों से ऊपर
शुद्ध वृहत् तृतीय; और $7$ ($3080/1760 = 1.75$) कोई कोमल लघु
सप्तम है --- अर्थात् विसंवादी।

\textbf{6.} $E = \tfrac14\mu L\omega^2A^2 = 0.25 \times 6.1 \times 10^{-3} \times 0.38 \times (2\pi
\times 440)^2 \times (5 \times 10^{-4})^2 = \qty{1.1}{mJ}$।

\textbf{7.} $A\sin kx\cos\omega t = \tfrac12A[\sin(kx - \omega t) + \sin(kx + \omega t)]$: अर्थात्
\qty{0.25}{mm} आयाम की दो तरंगें; $Z = \mu c = \qty{2.0}{kg/s}$; और हर एक
$\tfrac12Z\omega^2(A/2)^2 = \qty{0.49}{W}$ ले जाती है।

\textbf{8.} $\dot y(x, 0) = \sum_nB_n\omega_n\sin(n\pi x/L)$; इसे
$\sin(m\pi x/L)$ से गुणा कर समाकलित कीजिए, और $\int_0^L\sin\sin = \tfrac L2\delta_{mn}$ लीजिए।

\textbf{9.} समाकल्य $v_0\sin(n\pi x/L)$ है, $w$ चौड़ाई भर, $x_0$ के चारों ओर,
और वहाँ ज्या लगभग अचर रहती है।

\textbf{10.} $\sin(n\pi/8) = 0$ जब $n = 8, 16, \dots$; और चूँकि $B_n/B_1 =
\sin(n\pi/8)/(n\sin(\pi/8))$ है, सातवाँ संनादी घटकर पहले का $1/7$ रह जाता है,
जबकि संवादी $2$ से $6$ तक $0.9$ से $0.3$ बनाए रखते हैं: अर्थात् विसंवादी
सप्तम पालतू बना दिया जाता है।

\textbf{11.} प्रहार-बिंदु से दो वेग-स्पंद विपरीत दिशाओं में
निकलते हैं, सिरों पर उलटकर परावर्तित होते हैं, एक-दूसरे को पार करते और फिर जुड़ते हैं; और
आरंभिक अवस्था $2L/c = 1/f_1 = \qty{2.3}{ms}$ बाद लौट आती है।

\textbf{12.} $B = \pi^3 \times 2 \times 10^{11} \times 10^{-12}/(64 \times 685 \times 0.144) =
9.8 \times 10^{-4}$; $\sqrt{1 + 64B} = 1.031$: अर्थात् $+3.1\%$, और $1200\log_2(1.031) = 53$ सेंट।

\textbf{13.} किसी नीचे स्वर के संनादी ठीक गुणजों से तीव्र होते हैं;
और उन्हें ऊँचे स्वरों की मूल आवृत्तियों से मिलाने (तथा विस्पंद टालने) के लिए
ऊँचे स्वर कुछ तीव्र मिलाए जाते हैं।

\textbf{14.} $Z_s = \mu c = \qty{2.0}{kg/s}$; $r = (2 - 2000)/2002 = -0.998$।

\textbf{15.} $1 - r^2 \approx 4Z_s/Z_b = 4 \times 10^{-3}$।

\textbf{16.} प्रति सेकंड $c/2L = 440$ परावर्तन: अतः क्षय-दर $440 \times 4 \times
10^{-3} = \qty{1.8}{s^{-1}}$, $\tau = \qty{0.56}{s}$; और \qty{60}{dB} के लिए
$\ln10^6/1.8 = \qty{7.7}{s}$।

\textbf{17.} $1.1\,\text{mJ} \times 1.8\,\text{s}^{-1} = \qty{2}{mW}$ --- अर्थात् किसी धीमी बातचीत
से दो सौ गुना; और यही आप सुनते हैं।

\textbf{18.} \qty{1}{mm} का तार लगभग कोई हवा नहीं धकेलता: हवा उसके चारों ओर से
फिसल जाती है और उसकी विकिरित शक्ति नगण्य रहती है। ध्वनि-पट्ट
तार की ऊर्जा सेतु से लेकर वर्ग मीटर भर हवा हिलाता है:
अर्थात् क्षेत्रफल से प्रतिबाधा का मेल।

\textbf{19.} \qty{440}{Hz} A$_3$ का दूसरा संनादी है: अतः सेतु
उस तार को उसकी किसी विधा पर चला देता है --- यही सहानुभूति-अनुनाद है।

\textbf{20.} $Z_b$ बढ़ाइए: तब हर परावर्तन पर कम ऊर्जा रिसती है और स्वर अधिक
देर टिकता है --- पर वह धीमा भी हो जाता है।

\textbf{21.} $c = \qty{5.1}{km/s}$: पटरी में \qty{1.0}{s}, हवा में \qty{15}{s}।

\textbf{22.} $50/5064 = \qty{10}{ms}$; और मुक्त सिरे पर प्रतिबल का लुप्त होना
आवश्यक है: अतः संपीडन विरलन बनकर लौटता है।

\textbf{23.} $\lambda = \qty{2.5}{mm}$; $Z_{\text{इस्पात}} = \rho c = \qty{4e7}{kg.m^{-2}.s^{-1}}$
के सामने $400$: अतः $r \approx -1$, पूर्ण परावर्तन: कोई भी खुली दरार, चाहे कितनी
पतली हो, प्रतिध्वनि लौटा देती है।

\textbf{24.} $\Delta t = d(1/3.5 - 1/6) = d \times \qty{0.119}{s/km}$: $d = \qty{170}{km}$।

\textbf{25.} तार \qty{334}{m/s} ($T$, $\mu$); पटरी \qty{5.1}{km/s} ($E$, $\rho$);
चट्टान P \qty{6}{km/s} (संपीडन गुणांक, $\rho$); चट्टान S \qty{3.5}{km/s}
(अपरूपण गुणांक, $\rho$); हवा \qty{340}{m/s} ($\gamma P$, $\rho$ --- अगला अध्याय)।
\end{solution}
